\uuid{Bw9n}
\titre{Changement de variable, intégration par parties et intégrale impropre}
\niveau{L1}
\module{Analyse}
\chapitre{Calcul d'intégrales}
\sousChapitre{Changement de variables}
\theme{Intégration par parties, intégrale impropre, limite}
\auteur{}
\datecreate{2026-06-04}
\organisation{}
\difficulte{3}

\contenu{
\begin{enumerate}

\item
  \texte{On rappelle que $\sin^3 t = \sin t\,(1 - \cos^2 t)$.}
  \question{En considérant le changement de variable $u = \cos t$, calculer l'intégrale
    \[
      I = \int_0^\pi \frac{\sin^3 t}{1+\cos^2 t}\;dt.
    \]}
  \indication{Après le changement de variable, reconnaître que
    $\dfrac{1-u^2}{1+u^2} = \dfrac{2}{1+u^2} - 1$, et utiliser la parité de l'intégrande.}
  \reponse{
    On pose $u = \cos t$, d'où $du = -\sin t\,dt$.
    Pour $t = 0$ : $u = 1$ ; pour $t = \pi$ : $u = -1$.
    En utilisant $\sin^3 t\,dt = (1-\cos^2 t)(-du) = -(1-u^2)\,du$ :
    \[
      I = \int_1^{-1} \frac{1-u^2}{1+u^2}(-du) = \int_{-1}^1 \frac{1-u^2}{1+u^2}\,du.
    \]
    L'intégrande étant une fonction paire :
    \[
      I = 2\int_0^1 \frac{1-u^2}{1+u^2}\,du
        = 2\int_0^1\!\left(\frac{2}{1+u^2}-1\right)du
        = 2\Bigl[2\arctan u - u\Bigr]_0^1
        = 2\!\left(\frac{\pi}{2}-1\right)
        = \pi - 2.
    \]
  }

\item
  \question{Soit $\alpha > 2$. À l'aide d'une intégration par parties, calculer
    \[
      J(\alpha) = \int_2^\alpha \frac{\ln x}{x^2}\;dx.
    \]}
  \reponse{
    On pose $u(x) = \ln x$ et $v'(x) = x^{-2}$,
    d'où $u'(x) = \dfrac{1}{x}$ et $v(x) = -\dfrac{1}{x}$.
    \begin{align*}
      J(\alpha)
      &= \left[-\frac{\ln x}{x}\right]_2^\alpha
         + \int_2^\alpha \frac{1}{x} \cdot \frac{1}{x}\,dx \\
      &= -\frac{\ln\alpha}{\alpha} + \frac{\ln 2}{2}
         + \left[-\frac{1}{x}\right]_2^\alpha \\
      &= -\frac{\ln\alpha}{\alpha} + \frac{\ln 2}{2} - \frac{1}{\alpha} + \frac{1}{2} \\
      &= \frac{1+\ln 2}{2} - \frac{1+\ln\alpha}{\alpha}.
    \end{align*}
  }

\item
  \question{En déduire la valeur de $\displaystyle\int_2^{+\infty} \frac{\ln x}{x^2}\;dx$.}
  \reponse{
    Puisque $\ln\alpha = o(\alpha)$ quand $\alpha \to +\infty$, on a
    $\dfrac{1+\ln\alpha}{\alpha} \to 0$.
    Donc l'intégrale converge et
    \[
      \int_2^{+\infty} \frac{\ln x}{x^2}\,dx
      = \lim_{\alpha\to+\infty} J(\alpha)
      = \frac{1+\ln 2}{2}.
    \]
  }

\end{enumerate}
}
